\worldchapter{\trans{foes_title}}{appendix-foes}
\label{ch:appendix-foes}

\begin{multicols}{2}


    \section{Untoter}



    \subsection*{Vampirbrut}

    Die Vampirbrut ist ein niederer Vampir. Ohne großen Einfluss ist dies eine Gestalt, welche nur begrenzte Kräfte hat, und vielleicht auf sich selbst gestellt ist.

    \begin{tabular}{@{}ll@{}}
        \textbf{\trans{wounds_label}:} & 10 \\
        \textbf{\trans{movement_label}:} & 4 \\
        \textbf{\trans{strength_label}:} & 2 \\
        \textbf{\trans{dexterity_label}:} & 2 \\
        \textbf{\trans{mind_label}:} & 2 \\
        \textbf{\trans{resistances_label}:} & ['Kälte', 'Wasser'] \\
    \end{tabular}

    \textbf{Biss} (8)\\
    Bluten 1\\
    \textbf{Schwarze Magie} (6)\\
    Die Vampirbrut wirkt einen Zauber der Schwarzen Magie. Die Erfolge dieses Wurfs stellen die \textbf{Stärke} des Zaubers dar.\\


    \subsection*{Zombie}

    Der Zombie ist eine tragische Gestalt, welche in vielen Geschichten vorkommt. Ein Mensch, welcher auf besondere und unnatürliche Weise nach seinem Tod am Leben gehalten wird. Zombies haben ein Gehirn von der Größe einer Erbse, sie kennen nicht viel mehr als das Verlangen nach Blut. Und Gehirnen. Also, wenn man der Popkultur glauben darf.

    \begin{tabular}{@{}ll@{}}
        \textbf{\trans{wounds_label}:} & 6 \\
        \textbf{\trans{movement_label}:} & 3 \\
        \textbf{\trans{strength_label}:} & 3 \\
        \textbf{\trans{dexterity_label}:} & 1 \\
        \textbf{\trans{mind_label}:} & 1 \\
        \textbf{\trans{resistances_label}:} & ['Feuer', 'Gift'] \\
    \end{tabular}

    \textbf{Untoter Griff} (4)\\
    Vergiftet 1\\
    \textbf{Biss} (5)\\
    Vergiftet 1\\


    \subsection*{Mumie}

    Die lebendige Mumie ist ein Untoter. Sie ist auf magische Weise zum Leben erweckt worden. Nichts ist von ihrem Geist erhalten geblieben, alles, wonach sie trachtet, ist das Leben ihrer Opfer zu nehmen. Sie ist in der Regel unbewaffnet, ihr Fluch vergiftet jedoch ihre Opfer.

    \begin{tabular}{@{}ll@{}}
        \textbf{\trans{wounds_label}:} & 8 \\
        \textbf{\trans{movement_label}:} & 3 \\
        \textbf{\trans{strength_label}:} & 3 \\
        \textbf{\trans{dexterity_label}:} & 1 \\
        \textbf{\trans{mind_label}:} & 2 \\
        \textbf{\trans{resistances_label}:} & ['Feuer', 'Gift'] \\
    \end{tabular}

    \textbf{Griff} (6)\\
    Durchschlag 0\\
    \textbf{Fluch} (4)\\
    Durchschlag 3, Gift 2\\


    \subsection*{Skelett}

    Ein wandelndes Skelett, von dunkler Magie belebt.

    \begin{tabular}{@{}ll@{}}
        \textbf{\trans{wounds_label}:} & 4 \\
        \textbf{\trans{movement_label}:} & 3 \\
        \textbf{\trans{strength_label}:} & 2 \\
        \textbf{\trans{dexterity_label}:} & 2 \\
        \textbf{\trans{mind_label}:} & 2 \\
    \end{tabular}

    \textbf{Knochengriff} (8)\\
    Vergiftet 1\\
    \textbf{Kalter Atem} (6)\\
    Geschockt 2, 5 Meter Reichweite\\


    \section{Tier}



    \subsection*{Geier}

    Auf aufsteigenden Luftströmen kreist der Geier geduldig und beobachtet jede Regung unter sich. Er jagt nicht im klassischen Sinne, sondern wartet auf das Unvermeidliche, angezogen von Tod, Verfall und leichter Beute. Seine Präsenz wirkt wie ein stummes Omen und erscheint oft schon, bevor das Ende vollständig eingetreten ist. Am Boden bewegt er sich unbeholfen, doch mit unerbittlicher Effizienz beim Reißen und Fressen. In düsteren Auslegungen wirkt er wie ein Vorbote des Schicksals, der überall dort erscheint, wo Leben schwindet – als würde ihn das Sterben selbst rufen.

    \begin{tabular}{@{}ll@{}}
        \textbf{\trans{wounds_label}:} & 2 \\
        \textbf{\trans{movement_label}:} & 6 \\
        \textbf{\trans{strength_label}:} & 1 \\
        \textbf{\trans{dexterity_label}:} & 3 \\
        \textbf{\trans{mind_label}:} & 2 \\
    \end{tabular}

    \textbf{Schnabel} (7)\\
    Durchschlag 1\\


    \subsection*{Robbe}

    An kalten Küsten und treibendem Eis bewegt sich die Robbe mühelos zwischen Wasser und Land. Im Meer ist sie schnell und schwer fassbar, taucht und wendet sich mit fließender Präzision; an Land wirkt sie langsamer, bleibt jedoch wachsam und vorsichtig. Robben sind keine von Natur aus aggressiven Tiere, doch in die Enge getrieben oder beim Schutz ihres Nachwuchses reagieren sie mit plötzlicher Heftigkeit – beißen und schlagen mit überraschender Kraft um sich. Ihr Auftreten weist oft auf reiche Jagdgründe hin, aber auch auf verborgene Gefahren unter der Oberfläche. In düsteren Auslegungen wirken sie wachsam und wissend, ihre Augen knapp über der Wasserlinie, bevor sie lautlos verschwinden.

    \begin{tabular}{@{}ll@{}}
        \textbf{\trans{wounds_label}:} & 5 \\
        \textbf{\trans{movement_label}:} & 5 \\
        \textbf{\trans{strength_label}:} & 1 \\
        \textbf{\trans{dexterity_label}:} & 2 \\
        \textbf{\trans{mind_label}:} & 2 \\
    \end{tabular}

    \textbf{Biss} (5)\\
    Durchschlag 0\\


    \subsection*{Nashorn}

    Ein massiger und reizbarer Gigant, der für rohe Wucht, unbeugsame Zähigkeit und plötzliche Aggression steht. Sein schwaches Sehvermögen wird durch feines Gehör und Geruchssinn ausgeglichen, sodass es auf jede wahrgenommene Bedrohung sofort reagiert. Hat es sich einmal entschieden, stürmt es mit unaufhaltsamer Kraft voran und spießt oder zertrampelt alles auf seinem Weg. Ein Nashorn weicht nicht aus und zögert nicht – es durchbricht. In düsteren Auslegungen wirkt es wie ein lebender Rammbock, getrieben von blinder Wucht und unerschütterlichem Instinkt.

    \begin{tabular}{@{}ll@{}}
        \textbf{\trans{wounds_label}:} & 10 \\
        \textbf{\trans{movement_label}:} & 4 \\
        \textbf{\trans{strength_label}:} & 3 \\
        \textbf{\trans{dexterity_label}:} & 3 \\
        \textbf{\trans{mind_label}:} & 1 \\
    \end{tabular}

    \textbf{Ansturm} (10)\\
    Durchschlag 2\\


    \subsection*{Snake}

    Ein lautloser und schwer fassbarer Jäger, der für Geduld, Präzision und verborgene Tödlichkeit steht. Die Giftschlange verlässt sich auf Tarnung und regungslose Ruhe, um blitzschnell zuzuschlagen, sobald sich Beute nähert. Ihr Gift übernimmt den eigentlichen Tod – es schwächt, lähmt oder zerstört das Opfer langsam von innen. Schlangen meiden offene Konfrontationen und setzen auf Hinterhalt und Rückzug, doch ein einziger Fehler kann tödlich sein. In düsteren Auslegungen wirken sie kalt und berechnend, als sei jede Bewegung geplant und jeder Biss unausweichlich.

    \begin{tabular}{@{}ll@{}}
        \textbf{\trans{wounds_label}:} & 2 \\
        \textbf{\trans{movement_label}:} & 5 \\
        \textbf{\trans{strength_label}:} & 1 \\
        \textbf{\trans{dexterity_label}:} & 3 \\
        \textbf{\trans{mind_label}:} & 1 \\
    \end{tabular}

    \textbf{Biss} (8)\\
    Vergiftet 1\\


    \subsection*{Ratte}

    In Dunkelheit und verborgenen Winkeln gedeiht die Ratte dort, wo andere scheitern. Sie zwängt sich durch kleinste Spalten, nistet sich in Wänden ein und tritt oft in solcher Zahl auf, dass aus einer Plage eine Bedrohung wird. Allein ist sie scheu und leicht zu vertreiben, doch im Schwarm wird sie kühn – sie nagt, beißt und überwältigt durch schiere Hartnäckigkeit. Ratten verbreiten Schmutz, vernichten Vorräte und bringen unsichtbare Gefahren mit sich. In düsteren Auslegungen wirken sie wie eine schleichende Verseuchung, ein lebendiges Zeichen des Verfalls, das stärker wird, je länger man es ignoriert.

    \begin{tabular}{@{}ll@{}}
        \textbf{\trans{wounds_label}:} & 2 \\
        \textbf{\trans{movement_label}:} & 5 \\
        \textbf{\trans{strength_label}:} & 1 \\
        \textbf{\trans{dexterity_label}:} & 3 \\
        \textbf{\trans{mind_label}:} & 1 \\
    \end{tabular}

    \textbf{Biss} (5)\\
    Blutend 1\\


    \subsection*{Skorpion}

    Unter Sand, Stein oder im Schatten verborgen, verharrt der Skorpion in unheimlicher Regungslosigkeit. Wird er gestört, reagiert er augenblicklich – sein Stachel trifft mit präziser, geübter Geschwindigkeit. Seine Stärke liegt nicht in roher Gewalt, sondern im Gift, das schwächt, lähmt oder den Kampf sofort beendet. Er verfolgt nicht – er bestraft das Eindringen. In düsteren Auslegungen wirkt er uralt und gleichgültig, perfekt angepasst an eine Welt, in der Zögern den Tod bedeutet.

    \begin{tabular}{@{}ll@{}}
        \textbf{\trans{wounds_label}:} & 3 \\
        \textbf{\trans{movement_label}:} & 5 \\
        \textbf{\trans{strength_label}:} & 1 \\
        \textbf{\trans{dexterity_label}:} & 3 \\
        \textbf{\trans{mind_label}:} & 1 \\
    \end{tabular}

    \textbf{Biss} (6)\\
    Vergiftet 1\\
    \textbf{Stachel} (8)\\
    Vergiftet 2\\


    \subsection*{Frosch}

    Leuchtende Farben im Unterholz warnen vor Gefahr, lange bevor sich der Frosch selbst bewegt. Klein und leicht zu übersehen, verlässt er sich nicht auf Stärke, sondern auf hochwirksame Gifte, die selbst eine unachtsame Berührung gefährlich machen. Er jagt nicht durch Verfolgung, sondern durch Nähe – seine bloße Anwesenheit kann bereits schaden. In dichten Dschungeln oder feuchten Umgebungen wird er Teil der Landschaft, eine verborgene Gefahr zwischen Blättern und Wasser. In düsteren Auslegungen wirkt er beinahe unnatürlich, seine grelle Erscheinung wie eine bewusste Warnung, dass unter der Oberfläche etwas nicht stimmt.

    \begin{tabular}{@{}ll@{}}
        \textbf{\trans{wounds_label}:} & 2 \\
        \textbf{\trans{movement_label}:} & 4 \\
        \textbf{\trans{strength_label}:} & 1 \\
        \textbf{\trans{dexterity_label}:} & 1 \\
        \textbf{\trans{mind_label}:} & 1 \\
    \end{tabular}

    \textbf{Zunge} (3)\\
    Vergiftet 1\\


    \subsection*{Tiger}

    Ein einzelgängerischer und schwer fassbarer Prädator, der für Heimlichkeit, Geduld und plötzliche, explosive Gewalt steht. Anders als Rudeljäger verlässt er sich auf den perfekten Moment: ungesehen anschleichen, lautlos nähern und die Jagd mit einem einzigen, entscheidenden Angriff beenden. Tiger sind territorial und unberechenbar, verschwinden in ihrer Umgebung ebenso mühelos, wie sie daraus hervortreten. In düsteren Auslegungen wirken sie beinahe geisterhaft – lautlos, unausweichlich und tödlich präzise.

    \begin{tabular}{@{}ll@{}}
        \textbf{\trans{wounds_label}:} & 6 \\
        \textbf{\trans{movement_label}:} & 4 \\
        \textbf{\trans{strength_label}:} & 2 \\
        \textbf{\trans{dexterity_label}:} & 3 \\
        \textbf{\trans{mind_label}:} & 2 \\
    \end{tabular}

    \textbf{Biss} (10)\\
    Durchschlag 1, Blutend 1\\
    \textbf{Krallen} (10)\\
    Durchschlag 1\\


    \subsection*{Qualle}

    Lautlos und beinahe anmutig treibt die Qualle durch das Wasser und wirkt zunächst harmlos – bis es zur Berührung kommt. Ihr durchscheinender Körper verbirgt lange Tentakel, die mit starkem Gift versehen sind und augenblicklich reagieren. Sie jagt nicht im klassischen Sinne, sondern macht ihre Umgebung selbst zur Gefahr, in der jede Bewegung riskant wird. In Gruppen bilden sie fast unsichtbare Barrieren, die oft erst bemerkt werden, wenn es zu spät ist. In düsteren Auslegungen wirken sie wie lebende Fallen des Meeres – still, gleichgültig und unausweichlich.

    \begin{tabular}{@{}ll@{}}
        \textbf{\trans{wounds_label}:} & 2 \\
        \textbf{\trans{movement_label}:} & 4 \\
        \textbf{\trans{strength_label}:} & 1 \\
        \textbf{\trans{dexterity_label}:} & 2 \\
        \textbf{\trans{mind_label}:} & 1 \\
    \end{tabular}

    \textbf{Stachel} (4)\\
    Vergiftet 1\\


    \subsection*{Ochse}

    Ein kräftiges und eigensinniges Tier, das für rohe Gewalt, Ausdauer und territoriale Wut steht. Ob wild oder in Rage versetzt, reagiert es mit plötzlichen, heftigen Angriffen und setzt seine Hörner ein, um Gegner aufzuschlitzen oder fortzuschleudern. Anders als Raubtiere jagt es nicht – es stellt sich entgegen und vertreibt Eindringlinge durch schiere Wucht und Einschüchterung. Ochsen sind zwar kontrollierter, verfügen jedoch über enorme Kraft und können unter Stress oder Überlastung gefährlich werden. In düsteren Auslegungen wird der Bulle zum Sinnbild blinder Raserei und unaufhaltsamer Wucht – eine lebende Kraft, die sich nicht aufhalten lässt.

    \begin{tabular}{@{}ll@{}}
        \textbf{\trans{wounds_label}:} & 8 \\
        \textbf{\trans{movement_label}:} & 4 \\
        \textbf{\trans{strength_label}:} & 2 \\
        \textbf{\trans{dexterity_label}:} & 3 \\
        \textbf{\trans{mind_label}:} & 2 \\
    \end{tabular}

    \textbf{Rammen} (10)\\
    Durchschlag 1\\


    \subsection*{Elefant}

    Ein gewaltiger Riese der Natur, der für enorme Stärke, Widerstandskraft und unaufhaltsame Wucht steht. Meist ruhig und bedacht, wird er bei Bedrohung oder Provokation zu einer verheerenden Macht – er stürmt mit erdrückendem Gewicht voran und fegt alles beiseite, was sich ihm in den Weg stellt. Schon seine Größe verändert das Schlachtfeld und macht das Gelände zu seinem Vorteil. Elefanten sind nicht von Natur aus aggressiv, doch einmal erzürnt oder in Panik, sind sie kaum aufzuhalten. In düsteren Auslegungen wirken sie wie lebende Zerstörungsmaschinen – schwerfällig im Zorn, aber vernichtend in ihrer Wirkung.

    \begin{tabular}{@{}ll@{}}
        \textbf{\trans{wounds_label}:} & 10 \\
        \textbf{\trans{movement_label}:} & 4 \\
        \textbf{\trans{strength_label}:} & 4 \\
        \textbf{\trans{dexterity_label}:} & 2 \\
        \textbf{\trans{mind_label}:} & 2 \\
    \end{tabular}

    \textbf{Würgen} (8)\\
    Der Elefant würgt das Opfer mit seinem Rüssel.\\
    \textbf{Stampfen} (16)\\
    Das ganze Gewicht des Elefanten lastet auf dem Opfer. Der Angriff verursacht Durchschlag 2, kann aber nur ausgeführt werden ,wenn der Elefant bereits eine Runde neben dem Opfer steht.\\


    \subsection*{Adler}

    Hoch über dem Land zieht der Adler seine Kreise und überblickt mit scharfem Blick jede Bewegung unter sich. Er verschwendet keine Kraft – er wartet auf den richtigen Moment, bevor er in einem schnellen, tödlichen Sturzflug angreift. Seine Fänge treffen mit Präzision und Wucht und greifen die Beute, bevor sie reagieren kann. Selbst ohne Angriff beherrscht seine Präsenz den Himmel, ein ständiger Beobachter, dem man sich kaum entziehen kann. In düsteren Auslegungen wirkt er wie ein unerbittliches Auge des Himmels – fern, geduldig und absolut entschlossen, wenn er zuschlägt.

    \begin{tabular}{@{}ll@{}}
        \textbf{\trans{wounds_label}:} & 2 \\
        \textbf{\trans{movement_label}:} & 10 \\
        \textbf{\trans{strength_label}:} & 1 \\
        \textbf{\trans{dexterity_label}:} & 4 \\
        \textbf{\trans{mind_label}:} & 2 \\
    \end{tabular}

    \textbf{Krallen} (7)\\
    Blutend 1\\
    \textbf{Schnabel} (7)\\
    Durchschlag 1\\


    \subsection*{Fledermaus}

    Von Dachsparren und Höhlendecken stoßen Fledermäuse in rastloser Bewegung herab, geleitet eher von Klang als von Sicht. In unberechenbaren Schwärmen wirbeln sie durch die Luft, verwirren und beunruhigen, und ihre plötzliche Nähe kann selbst Erfahrene erschrecken. Allein sind sie zerbrechlich, doch in der Masse überwältigen sie – kratzen, beißen und erfüllen die Luft mit chaotischer Bewegung. Ihr Auftreten kündet oft von Verfall, Dunkelheit oder vergessenen Orten. In düsteren Auslegungen wirken sie wie lebendige Schatten, die aus der Finsternis hervorbrechen und ebenso schnell wieder verschwinden.

    \begin{tabular}{@{}ll@{}}
        \textbf{\trans{wounds_label}:} & 2 \\
        \textbf{\trans{movement_label}:} & 5 \\
        \textbf{\trans{strength_label}:} & 1 \\
        \textbf{\trans{dexterity_label}:} & 4 \\
        \textbf{\trans{mind_label}:} & 1 \\
    \end{tabular}

    \textbf{Biss} (2)\\
    Durchschlag 2\\


    \subsection*{Wildschwein}

    Eine brutale und furchtlose Naturgewalt, die für Aggression, Zähigkeit und territoriale Dominanz steht. Wird ein Wildschwein bedroht oder in die Enge getrieben, greift es mit explosiver Geschwindigkeit an und setzt seine Hauer ein, um Gegner aufzuschlitzen und niederzuwalzen. Wildschweine weichen kaum zurück – sie kämpfen kompromisslos weiter, selbst verletzt und im Schmerz. In düsteren Auslegungen wirken sie nahezu unaufhaltsam, getrieben von blinder Wut und dem unerbittlichen Instinkt, Eindringlinge zu vernichten.

    \begin{tabular}{@{}ll@{}}
        \textbf{\trans{wounds_label}:} & 8 \\
        \textbf{\trans{movement_label}:} & 4 \\
        \textbf{\trans{strength_label}:} & 3 \\
        \textbf{\trans{dexterity_label}:} & 1 \\
        \textbf{\trans{mind_label}:} & 2 \\
    \end{tabular}

    \textbf{Rammen} (10)\\
    Durchschlag 1\\


    \subsection*{Kamel}

    Ein widerstandsfähiges Tier rauer Landschaften, das in eine aggressive und gefährliche Form umschlägt. Kamele sind normalerweise zäh und bedächtig, können jedoch unter Stress oder in ihrem Revier äußerst unberechenbar werden und mit kräftigen Bissen, Tritten und plötzlichen Angriffen reagieren. Ihre Ausdauer erlaubt es ihnen, Gegner zu überdauern und Angriffe länger aufrechtzuerhalten, als man erwarten würde. Einmal erzürnt, lassen sie sich kaum einschüchtern – ihre Sturheit schlägt in unerbittliche Feindseligkeit um. In düsteren Auslegungen wirken sie unheimlich – zäh, nachtragend und beunruhigend hartnäckig, als würde die Wüste selbst sich wehren.

    \begin{tabular}{@{}ll@{}}
        \textbf{\trans{wounds_label}:} & 8 \\
        \textbf{\trans{movement_label}:} & 4 \\
        \textbf{\trans{strength_label}:} & 2 \\
        \textbf{\trans{dexterity_label}:} & 2 \\
        \textbf{\trans{mind_label}:} & 2 \\
    \end{tabular}

    \textbf{Biss} (6)\\
    Durchschlag 1\\


    \subsection*{Gorilla}

    Ein massiger und intelligenter Primat, der für kontrollierte Stärke, soziale Dominanz und plötzliche, explosive Gewalt steht. Innerhalb seiner Gruppe meist ruhig, reagiert er bei Bedrohung mit furchteinflößender Kraft und verteidigt Revier und Gefährten mit überwältigender physischer Gewalt. Dominanzgesten – Brusttrommeln, Brüllen, Scheinangriffe – dienen oft als Warnung, doch wenn er angreift, geschieht dies schnell und brutal. Ein Gorilla jagt nicht wie ein Raubtier; er stellt sich der Gefahr und zerschmettert sie. In düsteren Auslegungen wirkt er beinahe bewusst in seinem Zorn – berechnend, beschützend und unaufhaltsam, wenn er entfesselt wird.

    \begin{tabular}{@{}ll@{}}
        \textbf{\trans{wounds_label}:} & 6 \\
        \textbf{\trans{movement_label}:} & 4 \\
        \textbf{\trans{strength_label}:} & 2 \\
        \textbf{\trans{dexterity_label}:} & 3 \\
        \textbf{\trans{mind_label}:} & 3 \\
    \end{tabular}

    \textbf{Griff} (8)\\
    Das Opfer wird festgehalten und kann sich nur mit einer Kraft Probe befreien.\\


    \subsection*{Eule}

    In den stillen Stunden bewegt sich die Eule dort, wo selbst Geräusche zu verschwinden scheinen. Ihr Flug ist nahezu lautlos, die Schwingen gleiten ohne Vorwarnung durch die Dunkelheit. Geleitet von außergewöhnlichem Gehör nimmt sie Bewegungen wahr, die anderen entgehen, und schlägt mit plötzlicher Präzision von oben zu. Sie verweilt nicht – erscheint nur kurz, um zu greifen, und verschwindet ebenso schnell wieder im Schatten. In düsteren Auslegungen wirkt sie wie eine Beobachterin der Nacht, ungesehen und doch stets präsent, ein Flüstern von Bewegung jenseits der Wahrnehmung.

    \begin{tabular}{@{}ll@{}}
        \textbf{\trans{wounds_label}:} & 2 \\
        \textbf{\trans{movement_label}:} & 8 \\
        \textbf{\trans{strength_label}:} & 1 \\
        \textbf{\trans{dexterity_label}:} & 3 \\
        \textbf{\trans{mind_label}:} & 4 \\
    \end{tabular}

    \textbf{Schnabel} (6)\\
    Durchschlag 0\\


    \subsection*{Gans}

    An stillen Gewässern und auf offenen Flächen behauptet die Gans ihr Revier mit unerwarteter Entschlossenheit. Was zunächst harmlos wirkt, wird bei Annäherung schnell konfrontativ – sie zischt, breitet die Flügel aus und geht ohne Zögern nach vorn. Mit überraschender Kühnheit verteidigt sie ihr Gebiet und ihren Nachwuchs, hackt mit dem Schnabel zu und schlägt mit kräftigen Flügeln. Rückzug ist nicht ihr erster Impuls – sondern Angriff. In düsteren Auslegungen wirkt sie fast unerbittlich, eine kleine, aber wütende Wächterin, die sich unter keinen Umständen vertreiben lässt.

    \begin{tabular}{@{}ll@{}}
        \textbf{\trans{wounds_label}:} & 2 \\
        \textbf{\trans{movement_label}:} & 6 \\
        \textbf{\trans{strength_label}:} & 1 \\
        \textbf{\trans{dexterity_label}:} & 2 \\
        \textbf{\trans{mind_label}:} & 3 \\
    \end{tabular}

    \textbf{Biss} (8)\\
    Blutend 1\\


    \subsection*{Wilder Hund}

    Ein unerbittlicher Rudeljäger, der von Koordination, Ausdauer und Instinkt lebt. Wilde Hunde verlassen sich nicht auf reine Stärke – sie siegen durch Zusammenarbeit, Kommunikation und das Zermürben ihrer Beute. Über lange Strecken treiben sie ihr Ziel vor sich her, erzwingen Fehler und isolieren Schwache, bevor sie gemeinsam zuschlagen. Einzelne sind beherrschbar, doch als Rudel werden sie zu einer schnellen, disziplinierten Bedrohung. In düsteren Auslegungen wirken sie unheimlich synchron – wie ein einziger Wille, verteilt auf viele Körper.

    \begin{tabular}{@{}ll@{}}
        \textbf{\trans{wounds_label}:} & 6 \\
        \textbf{\trans{movement_label}:} & 4 \\
        \textbf{\trans{strength_label}:} & 2 \\
        \textbf{\trans{dexterity_label}:} & 2 \\
        \textbf{\trans{mind_label}:} & 2 \\
    \end{tabular}

    \textbf{Biss} (6)\\
    Blutend 1\\
    \textbf{Ruf nach dem Rudel} (6)\\
    Der wilde Hund ruft nach seinem Rudel. Nach 1W3 Kampfrunden erscheinen weitere wilde Hunde entsprechend der Erfolge dieses Wurfs.\\


    \subsection*{Riesenkrabbe}

    Über Gezeitenflächen und entlang zerklüfteter Küsten bewegt sich die Riesenkrabbe mit beunruhigender Zielstrebigkeit. Ihr schwerer Panzer hält selbst kräftige Schläge ab, während ihre gewaltigen Scheren mit zermalmender Kraft zuschnappen und Knochen wie Rüstung brechen können. Sie stürmt nicht blind voran – jede Bewegung ist bedacht, während sie sich nähert und sich zugleich mit erhobenen Scheren schützt. In ihrem Revier ist der Untergrund tückisch, und die Flucht wird durch das Gelände erschwert, das sie bestens kennt. In düsteren Auslegungen wirkt sie wie eine lebende Bastion der Küste – unnachgiebig, fremdartig und darauf ausgelegt, Unvorsichtige einem langsamen, unausweichlichen Ende entgegenzuführen.

    \begin{tabular}{@{}ll@{}}
        \textbf{\trans{wounds_label}:} & 6 \\
        \textbf{\trans{movement_label}:} & 6 \\
        \textbf{\trans{strength_label}:} & 2 \\
        \textbf{\trans{dexterity_label}:} & 2 \\
        \textbf{\trans{mind_label}:} & 1 \\
    \end{tabular}

    \textbf{Klauen} (10)\\
    \textit{Dad-a-chum? Dum-a-chum? Ded-a-chek? Did-a-chick?}\\


    \subsection*{Wolf}

    Ein räuberischer Jäger der Wildnis, der für Geduld, Instinkt und Zusammenarbeit steht. Wölfe treten selten allein auf, sondern jagen im Rudel, prüfen ihre Beute auf Schwäche und schlagen erst dann zu. Sie meiden unnötige Risiken, kreisen ihr Ziel ein, zermürben es und greifen an, wenn der Vorteil sicher ist. In düsteren Auslegungen wirken sie unnatürlich klug, unerbittlich oder beunruhigend furchtlos.

    \begin{tabular}{@{}ll@{}}
        \textbf{\trans{wounds_label}:} & 4 \\
        \textbf{\trans{movement_label}:} & 5 \\
        \textbf{\trans{strength_label}:} & 2 \\
        \textbf{\trans{dexterity_label}:} & 3 \\
        \textbf{\trans{mind_label}:} & 2 \\
    \end{tabular}

    \textbf{Biss} (7)\\
    Blutend 1\\


    \subsection*{Spinne}

    Ein geduldiger und methodischer Jäger, der für Regungslosigkeit, Präzision und das Einfangen steht. Die giftige Spinne wartet verborgen in Netzen oder dunklen Winkeln und schlägt nur zu, wenn ihre Beute gefangen oder wehrlos ist. Ihr Gift schwächt oder lähmt, sodass ein Entkommen kaum möglich ist. Anders als verfolgenden Jägern gehört ihr die Umgebung selbst – sie macht den Raum zur Falle. In düsteren Auslegungen wirkt sie berechnend und fremdartig, eine lautlose Architektin des Todes, deren Präsenz oft erst bemerkt wird, wenn es zu spät ist.

    \begin{tabular}{@{}ll@{}}
        \textbf{\trans{wounds_label}:} & 2 \\
        \textbf{\trans{movement_label}:} & 2 \\
        \textbf{\trans{strength_label}:} & 1 \\
        \textbf{\trans{dexterity_label}:} & 3 \\
        \textbf{\trans{mind_label}:} & 1 \\
    \end{tabular}

    \textbf{Biss} (4)\\
    Vergiftet 1, Durchschlag 1\\


    \subsection*{Riesenkraken}

    Aus der Tiefe entfaltet sich etwas Großes und Geduldiges. Der Riesenkalmar ist ein Meister der Tarnung, verschmilzt mit Fels und Riff, bis er sich entscheidet zu handeln. Seine Intelligenz wirkt beunruhigend – prüfend, tastend, lernend. Wenn er zuschlägt, dann mit plötzlicher Reichweite, umschlingt seine Beute mit kraftvollen Armen und zieht sie in einen erdrückenden, unausweichlichen Griff. Das Wasser selbst wird zur Waffe, verschleiert Sicht und Orientierung. In düsteren Auslegungen wirkt er uralt und wachsam, ein lauernder Verstand in der Tiefe, der auf den richtigen Moment wartet.

    \begin{tabular}{@{}ll@{}}
        \textbf{\trans{wounds_label}:} & 10 \\
        \textbf{\trans{movement_label}:} & 4 \\
        \textbf{\trans{strength_label}:} & 3 \\
        \textbf{\trans{dexterity_label}:} & 3 \\
        \textbf{\trans{mind_label}:} & 2 \\
    \end{tabular}

    \textbf{Würgegriff} (10)\\
    Das Opfer ist gegriffen und erleidet jede Kampfrunde 2 Wunden. Es kann sich in seiner Kampfrunde mit einer Kraft Probe befreien.\\


    \subsection*{Bär}

    Eine gewaltige und widerstandsfähige Naturgewalt, die für rohe Stärke, Ausdauer und territoriale Wut steht. Anders als schnelle Jäger verlässt sich der Bär nicht auf Geschwindigkeit oder Heimlichkeit – einmal aufgebracht, überwältigt er seine Gegner durch schiere Kraft und unerbittliche Härte. Meist wirkt er gleichgültig, doch bei Bedrohung oder zur Verteidigung seines Reviers wird er zu einer verheerenden Macht. In düsteren Auslegungen erscheint er nahezu unaufhaltsam – die Verkörperung urtümlicher Raserei.

    \begin{tabular}{@{}ll@{}}
        \textbf{\trans{wounds_label}:} & 8 \\
        \textbf{\trans{movement_label}:} & 3 \\
        \textbf{\trans{strength_label}:} & 3 \\
        \textbf{\trans{dexterity_label}:} & 2 \\
        \textbf{\trans{mind_label}:} & 2 \\
    \end{tabular}

    \textbf{Biss} (8)\\
    Blutend 2\\
    \textbf{Krallen} (8)\\
    Blutend 1\\


    \subsection*{Krokodil}

    Ein geduldiger, urtümlicher Lauerjäger, der für Regungslosigkeit, Präzision und plötzliche, überwältigende Gewalt steht. Halb im Wasser verborgen, nahezu unsichtbar, schlägt er in einem Augenblick zu, wenn Beute dem Ufer zu nahe kommt. Sein Griff ist unerbittlich – einmal gepackt, gibt es kaum ein Entkommen. Krokodile verfolgen ihre Beute nicht weit; sie beherrschen die Grenze zwischen Land und Wasser und verwandeln sicheren Boden in eine tödliche Falle. In düsteren Auslegungen wirken sie uralt und berechnend, als hätte das Wasser selbst Zähne.

    \begin{tabular}{@{}ll@{}}
        \textbf{\trans{wounds_label}:} & 6 \\
        \textbf{\trans{movement_label}:} & 4 \\
        \textbf{\trans{strength_label}:} & 2 \\
        \textbf{\trans{dexterity_label}:} & 2 \\
        \textbf{\trans{mind_label}:} & 2 \\
    \end{tabular}

    \textbf{Biss} (10)\\
    Das Opfer kann sich nur mit einer Kraft Probe aus dem Biss befreien. Das Krokodil zieht gerne Opfer unter Wasser, wenn dies möglich ist.\\


    \subsection*{Hai}

    Unter der Oberfläche bewegt sich der Hai mit lautloser Sicherheit und nimmt wahr, was unsichtbar bleibt. Er stürzt sich nicht blind auf seine Beute – er kreist, prüft und nähert sich mit kalkulierter Entschlossenheit. Ein plötzlicher Angriff beendet die Begegnung in einem Augenblick aus Bewegung und Blut, bevor er wieder in der Tiefe verschwindet. In seinem Element gibt es keinen festen Halt und kaum ein Entkommen. In düsteren Auslegungen wirkt er wie eine unsichtbare Macht der Tiefe, angezogen von Schwäche und Unruhe – als hätte sich das Wasser selbst gegen seine Beute gewandt.

    \begin{tabular}{@{}ll@{}}
        \textbf{\trans{wounds_label}:} & 8 \\
        \textbf{\trans{movement_label}:} & 6 \\
        \textbf{\trans{strength_label}:} & 2 \\
        \textbf{\trans{dexterity_label}:} & 3 \\
        \textbf{\trans{mind_label}:} & 1 \\
    \end{tabular}

    \textbf{Biss} (10)\\
    Durchschlag 1\\


    \subsection*{Hyene}

    Ein listiger Aasfresser und Gelegenheitsjäger, der dort überlebt, wo andere scheitern. Die Hyäne lebt von Anpassungsfähigkeit, Ausdauer und rücksichtsloser Zweckmäßigkeit und meidet oft den direkten Angriff, bis sich eine Schwäche zeigt. Sie prüfen ihre Beute unermüdlich, zermürben sie und nutzen jede Gelegenheit aus. Ihre unheimlichen Laute und ihr verstörendes Auftreten verleihen ihnen einen beinahe widernatürlichen Ruf. In düsteren Auslegungen wirken sie spöttisch und grausam – Kreaturen, die sich am Verfall und am Untergang anderer laben.

    \begin{tabular}{@{}ll@{}}
        \textbf{\trans{wounds_label}:} & 4 \\
        \textbf{\trans{movement_label}:} & 4 \\
        \textbf{\trans{strength_label}:} & 1 \\
        \textbf{\trans{dexterity_label}:} & 3 \\
        \textbf{\trans{mind_label}:} & 2 \\
    \end{tabular}

    \textbf{Biss} (8)\\
    Blutend 1\\


    \subsection*{Löwe}

    Ein dominanter Jäger, der für Stärke, Autorität und tödliche Präzision steht. Anders als ausdauernde Rudeljäger verlässt er sich auf kurze, überwältigende Angriffe, oft aus dem Hinterhalt, oder behauptet sein Revier allein durch seine Präsenz. Ein Löwe jagt nicht endlos – er entscheidet, wann die Jagd beginnt und endet. In düsteren Auslegungen wirkt er königlich und furchteinflößend, ein Sinnbild von Macht, das keinen Widerspruch duldet.

    \begin{tabular}{@{}ll@{}}
        \textbf{\trans{wounds_label}:} & 6 \\
        \textbf{\trans{movement_label}:} & 4 \\
        \textbf{\trans{strength_label}:} & 3 \\
        \textbf{\trans{dexterity_label}:} & 2 \\
        \textbf{\trans{mind_label}:} & 2 \\
    \end{tabular}

    \textbf{Biss} (8)\\
    Durchschlag 1, Blutend 1\\
    \textbf{Krallen} (6)\\
    Piercing 1, Bleeding 2\\


    \section{Magisch}



    \subsection*{Werwolf}

    Der Werwolf ist ein Lebewesen, halb Wolf, halb Mensch. Bei Vollmond verwandelt sich der bis dahin normale Mensch in einen Wolf, welcher jedoch bis zu zwei Schritt groß werden kann und äußerst aggressiv ist. Jeder kann ein Werwolf werden, sobald er von einem bereits infizierten Werfolf gebissen wird.

    \begin{tabular}{@{}ll@{}}
        \textbf{\trans{wounds_label}:} & 12 \\
        \textbf{\trans{movement_label}:} & 8 \\
        \textbf{\trans{strength_label}:} & 5 \\
        \textbf{\trans{dexterity_label}:} & 3 \\
        \textbf{\trans{mind_label}:} & 3 \\
    \end{tabular}

    \textbf{Biss} (8)\\
    Durchschlag 1, Blutend 1\\
    \textbf{Krallen} (8)\\
    Blutend 1\\


    \subsection*{Kompostfee}

    Die Kompostfee ist ein sonderbares, magisches Wesen. Sie lebte einst im Kompost der Hexe Mare, im Mittelalter auf der Erde, unweit der Stadt Aquisgrani. Mare und sie verbinden ein besonderes Band, welches auch auf magische Weise geprägt ist.

    \begin{tabular}{@{}ll@{}}
        \textbf{\trans{wounds_label}:} & 2 \\
        \textbf{\trans{movement_label}:} & 8 \\
        \textbf{\trans{strength_label}:} & 2 \\
        \textbf{\trans{dexterity_label}:} & 5 \\
        \textbf{\trans{mind_label}:} & 3 \\
        \textbf{\trans{resistances_label}:} & ['Magie'] \\
    \end{tabular}

    \textbf{Energiekugel} (2)\\
    Durchschlag 2, Geschockt 1\\


    \subsection*{Schillerwal (Leviatan)}

    Der Schillerwal ähnelt in seiner Form einem irdischen Blauwal, ist jedoch schlanker und besitzt längere, fast flügelartige Seitenflossen. Das Besondere ist seine Haut: Sie ist nicht grau, sondern besitzt eine perlmuttartige, semitransparente Oberfläche.

Je nach Lichteinfall und magischer Sättigung der Umgebung bricht seine Haut das Licht in alle Farben des Spektrums. Daher der Name. Wenn ein Schillerwal an der Oberfläche bricht, sieht es aus, als würde ein lebender Regenbogen aus dem Wasser steigen.
Gelehrte glauben, dass diese Wale sich nicht nur von Krill ernähren, sondern auch das Licht des Mondes und der Sterne absorbieren, wenn sie nachts an die Oberfläche kommen.

Schillerwale meiden die flachen Küstengewässer. Sie ziehen durch die tiefen Ozeane, weit abseits der Routen der Handelsschiffe.

Man sichtet sie am häufigsten im Südmeer, in den kalten Strömungen fernab der Hitze des Dschungels, oder in den mystischen Gewässern rund um verlassene Inselarchive.

Sie reisen in kleinen Familienverbänden (Pods). Es heißt, ihr Gesang kann Stürme beruhigen oder Wahnsinn bei jenen auslösen, die ihn zu lange hören.

Da die Jagd extrem gefährlich ist (Schillerwale wehren sich selten, werden aber oft von See-Elementaren oder Meerjungfrauen beschützt) und die Tiere selten sind, ist der Preis enorm.

Die Ancatir betrachten die Jagd auf Schillerwale als Frevel. Sie nutzen nur Tran, der von natürlich gestrandeten Walen stammt (Geschenk der Gezeiten). Alchemisten, die blutigen Tran verwenden, werden in elfischen Enklaven oft der Stadt verwiesen.

Der Schillerwal gehört zu den faszinierendsten und friedliebendsten Giganten der Meere um Tirakan. Er ist nicht nur ein Tier, sondern eine lebende Anomalie, die eng mit den magischen Strömungen der Ozeane verbunden ist.

Schillerwale greifen niemals zuerst an. Wird er getötet, erleiden alle Beteiligten einen Fluch des Meeres (+2 Erschöpfung pro Tag auf See).

Wir sahen ihn bei Neumond. Er leuchtete unter dem Kiel wie eine versunkene Stadt. Als wir die Harpunen warfen, schrie das Biest nicht. Es fing an zu singen. Ein Ton, so tief, dass das Holz meines Schiffs splitterte und zwei meiner Männer einfach ins Wasser sprangen, mit einem Lächeln im Gesicht. Wir haben ihn erlegt, ja. Aber das Öl... es brennt in den Lampen meiner Kajüte, und ich schwöre, in den Schatten sehe ich die Gesichter derer, die gesprungen sind.
Aus dem Logbuch des Walfängers 'Haken-Ulf'

    \begin{tabular}{@{}ll@{}}
        \textbf{\trans{wounds_label}:} & 80 \\
        \textbf{\trans{movement_label}:} & 20 \\
        \textbf{\trans{strength_label}:} & 16 \\
        \textbf{\trans{dexterity_label}:} & 12 \\
        \textbf{\trans{mind_label}:} & 11 \\
    \end{tabular}

    \textbf{Sonarstoß} (16)\\
    Der Wal stößt einen extrem lauten, gerichteten Schallimpuls aus.

Die Ziele müssen eine Resistenz-Probe (WI) bestehen. Bei Misserfolg sind sie für 1W6 Runden betäubt und erleiden einen Malus von +4 auf alle Proben.\\
    \textbf{Gewaltiger Flossenschlag} (14)\\
    Ein Schlag mit der massiven Schwanzflosse, der das Wasser peitscht.

3W6+16 Schaden. Trifft alle Ziele in einem 5-Meter-Radius hinter dem Wal.\\
    \textbf{Irisierendes Leuchten} (14)\\
    Der Wal lässt seine Haut in einem blendenden Farbspektrum aufblitzen.

Alle Angreifer im Nahbereich müssen eine Wahrnehmungs-Probe bestehen. Bei Misserfolg sind sie für die nächste Runde geblendet und können keine Angriffe ausführen.\\
    \textbf{Abtauchen} (12)\\
    Der Wal zieht sich in die Tiefe zurück.

Erhöht den RS für eine Runde um weitere +4 Schilde, während er sich aus dem Kampfbereich entfernt.\\

\end{multicols}